\documentclass[11pt]{article}

% ---------- Packages ----------
\usepackage[utf8]{inputenc}
\usepackage[margin=1in]{geometry}
\usepackage{graphicx}
\usepackage{amsmath}
\usepackage{amssymb}
\usepackage{listings}
\usepackage{xcolor}
\usepackage{caption}
\usepackage{subcaption}
\usepackage{hyperref}

% ---------- Code listing style (for pasted code, optional) ----------
\definecolor{codegray}{rgb}{0.5,0.5,0.5}
\definecolor{codegreen}{rgb}{0,0.5,0}
\definecolor{codeblue}{rgb}{0.1,0.1,0.7}
\lstdefinestyle{py}{
    language=Python,
    basicstyle=\ttfamily\footnotesize,
    keywordstyle=\color{codeblue},
    commentstyle=\color{codegray}\itshape,
    stringstyle=\color{codegreen},
    numbers=left,
    numberstyle=\tiny\color{codegray},
    breaklines=true,
    frame=single,
    showstringspaces=false,
    tabsize=4
}

\title{Problem 2: Sphere Tracing and Volumetric Rendering}
\author{}
\date{}

\begin{document}
\maketitle

%==================================================================
\section{Problem 2.1: Sphere Tracing}
%==================================================================

\subsection{Overview}
In this part I implemented a sphere tracing renderer for an implicit scene made of two
spheres (blue and green), each defined by its Signed Distance Function (SDF). The pipeline
has two stages: (i) generating one ray per pixel from the camera parameters, and (ii)
marching each ray forward by the SDF value until it reaches a surface.

\subsection{Code}
% -------- Paste / screenshot your code here --------
% You can either screenshot the functions and use \includegraphics,
% or paste the source directly using the listings block below.

% --- Option A: screenshot ---
% \begin{figure}[h]
%     \centering
%     \includegraphics[width=\linewidth]{images/code_sphere_tracing.png}
%     \caption{\texttt{camera\_param\_to\_rays} and \texttt{sphere\_tracing}.}
% \end{figure}

% --- Option B: paste the code ---
% \begin{lstlisting}[style=py]
% # paste camera_param_to_rays + sphere_tracing here
% \end{lstlisting}

\vspace{1cm}
\textit{[ Paste the code for \texttt{camera\_param\_to\_rays} and \texttt{sphere\_tracing} here ]}
\vspace{1cm}

\subsection{Rendered Result (\texttt{submit} intrinsics)}
% -------- Place the sphere tracing render here --------
\begin{figure}[h]
    \centering
    % \includegraphics[width=0.9\linewidth]{images/sphere_tracing.png}
    \fbox{\parbox[c][5cm][c]{0.9\linewidth}{\centering \textit{[ sphere\_tracing.png goes here ]}}}
    \caption{Sphere tracing result for the two camera views (blue and green spheres).}
    \label{fig:sphere_tracing}
\end{figure}

%==================================================================
\section{Problem 2.2: Differentiable Volume Rendering}
%==================================================================

\subsection{Overview}
Here I implemented a differentiable volume renderer. Instead of jumping to the surface, the
ray is split into \texttt{num\_samples} sample points between $t_{near}$ and $t_{far}$. At
every sample the implicit model returns a density $\sigma_i$ and a color $c_i$, and the final
pixel color is computed with the classical volume rendering equation
\begin{equation}
    C = \sum_i T_i\,\alpha_i\,c_i, \qquad
    \alpha_i = 1 - e^{-\sigma_i \delta_i}, \qquad
    T_i = \prod_{j=1}^{i-1}\left(1-\alpha_j\right),
\end{equation}
where $\delta_i$ is the distance between adjacent samples and $T_i$ is the accumulated
transmittance (how much light survives up to sample $i$).

\subsection{Code}
\vspace{1cm}
\textit{[ Paste the code for \texttt{camera\_param\_to\_rays}, \texttt{sample\_points\_on\_rays} and \texttt{volume\_rendering} here ]}
\vspace{1cm}

\subsection{Rendered Result (\texttt{submit} intrinsics)}
% -------- Place the volume rendering render here --------
\begin{figure}[h]
    \centering
    % \includegraphics[width=0.9\linewidth]{images/volume_rendering.png}
    \fbox{\parbox[c][5cm][c]{0.9\linewidth}{\centering \textit{[ volume\_rendering.png goes here ]}}}
    \caption{Volume rendering result for the two camera views
    (blue/green spheres, red/yellow cubes).}
    \label{fig:volume_rendering}
\end{figure}

%==================================================================
\subsection{Question: Why is volume rendering preferred over sphere tracing?}
%==================================================================

In modern differentiable rendering, volumetric rendering is used far more often than sphere
tracing. There are a few reasons for this:

\begin{itemize}
    \item \textbf{It is smooth and differentiable everywhere.} Volume rendering integrates a
    continuous density field, so the final color is a smooth function of the scene parameters
    and \emph{every} sample along the ray receives a gradient. Sphere tracing, on the other
    hand, returns a single hard surface hit and uses a variable number of marching steps with
    a hard threshold ($\text{SDF} < \epsilon$). This makes its gradients sparse and poorly
    defined, which is bad for optimization.

    \item \textbf{Not everything in the world is a well-defined surface.} Sphere tracing
    fundamentally assumes the scene is a solid surface described by a valid SDF. But many
    real phenomena -- smoke, fog, clouds, fur, hair, fire, or semi-transparent materials --
    simply do not have a clean surface. A density field naturally represents these soft,
    volumetric, and semi-transparent effects, while an SDF cannot.

    \item \textbf{It does not require a valid SDF.} Volume rendering works with an arbitrary
    density field, so a network is free to learn any volume without being forced to satisfy
    the strict distance (Eikonal) constraint that an SDF must obey. This makes it much easier
    to learn a scene from images, which is exactly what methods like NeRF rely on.

    \item \textbf{Soft accumulation handles edges and uncertainty.} Because the color is a
    weighted average over many samples, thin structures, object boundaries, and uncertain
    regions all contribute gradients, instead of the optimization depending on one surface
    point. This generally leads to more stable training.
\end{itemize}

In short: volume rendering gives clean gradients to many samples and can represent both solid
objects and fuzzy volumetric stuff, which is why it became the standard tool for
differentiable and neural rendering.

%==================================================================
\section{Extra Challenge}
%==================================================================

\subsection{Effect of the number of samples}
% -------- Place renders with different sample counts here --------
\begin{figure}[h]
    \centering
    % \includegraphics[width=0.9\linewidth]{images/samples_comparison.png}
    \fbox{\parbox[c][5cm][c]{0.9\linewidth}{\centering \textit{[ renders with different num\_samples go here ]}}}
    \caption{Effect of changing the number of samples along each ray
    (e.g.\ 16 / 64 / 256 samples).}
    \label{fig:samples}
\end{figure}

\noindent\textit{[ Short discussion: how does rendering quality change as the number of
samples increases? ]}

\subsection{Stratified sampling}
% -------- Place stratified sampling visualization here --------
\begin{figure}[h]
    \centering
    % \includegraphics[width=0.9\linewidth]{images/stratified_sampling.png}
    \fbox{\parbox[c][5cm][c]{0.9\linewidth}{\centering \textit{[ stratified sampling visualization goes here ]}}}
    \caption{Uniform sampling vs.\ stratified sampling.}
    \label{fig:stratified}
\end{figure}

\noindent\textit{[ Short discussion: implementation of stratified sampling and the difference
you observe compared to uniform sampling. ]}

\end{document}
